% glossary.tex -- Module 8 glossary (A-Z). GENERATED by
% script) from the \term / \kw words of the
\unnumberedchapter[\faBook]{Glossary}

Every technical word in this book, from A to Z. The meaning is in plain English. The
last column gives the usual Chinese term used in physics textbooks.

\begin{glossarytable}
  % ---- A
  \gterm{Absolute magnitude ($M$)}{brightness at \qty{10}{pc}; smaller = brighter}{绝对星等}
  \gterm{Absorption spectrum}{dark lines on a continuous spectrum}{吸收光谱}
  \gterm{Accurate (accuracy)}{close to the true or accepted value}{准确度}
  \gterm{Activity ($A$)}{decays per second, in becquerel (Bq)}{放射性活度}
  \gterm{Alpha particle ($\alpha$)}{a helium-4 nucleus, charge $+2e$}{α粒子}
  \gterm{Angular momentum}{amount of turning motion; $mvr$ for an orbit}{角动量}
  \gterm{Annealing}{heating a metal, then cooling it slowly}{退火}
  \gterm{Annihilation}{particle + antiparticle turn into photons}{湮灭}
  \gterm{Anode}{the positive electrode of a tube}{阳极}
  \gterm{Antineutrino}{antiparticle of the neutrino (in $\beta^-$ decay)}{反中微子}
  \gterm{Antiparticle}{same mass; opposite charge and properties}{反粒子}
  \gterm{Apparent brightness}{how bright a star looks from Earth}{视亮度}
  \gterm{Artificial transmutation}{a nuclear change forced by firing a particle}{人工核转变}
  \gterm{Atomic mass}{mass of a whole atom, electrons included}{原子质量}
  \gterm{Atomic mass unit (u)}{\qty{1}{\amu} $=$ \qty{1.661e-27}{kg} $=$ \qty{931.5}{MeV}/$c^2$}{原子质量单位}
  \gterm{Atomic number ($Z$)}{number of protons; it decides the element}{原子序数}
  \gterm{Atomiser}{a sprayer that makes a fine mist of drops}{喷雾器}
  % ---- B
  \gterm{Background radiation}{radiation that is always around us}{本底辐射}
  \gterm{Balmer series}{hydrogen lines ending on $n = 2$ (four visible)}{巴耳末系}
  \gterm{Band of stability}{the strip of stable nuclei on an $N$--$Z$ graph}{稳定带}
  \gterm{Baryon}{a hadron of three quarks, such as the proton}{重子}
  \gterm{Becquerel (Bq)}{unit of activity: one decay per second}{贝可勒尔}
  \gterm{Beta particle ($\beta$)}{a fast electron ($\beta^-$) or positron ($\beta^+$)}{β粒子}
  \gterm{Big Bang}{the hot, dense start of the Universe}{大爆炸}
  \gterm{Big Bang nucleosynthesis}{helium nuclei made 3--20 minutes after it}{大爆炸核合成}
  \gterm{Binding energy ($E_\mathrm{B}$)}{energy needed to pull a nucleus apart}{结合能}
  \gterm{Binding energy per nucleon}{$E_\mathrm{B} \div A$; bigger = more tightly bound}{比结合能}
  \gterm{Black body}{ideal absorber; its glow depends only on $T$}{黑体}
  \gterm{Black hole}{a collapsed core that even light cannot leave}{黑洞}
  \gterm{Blueshift}{lines moved to shorter wavelengths}{蓝移}
  \gterm{Bohr radius}{radius of Bohr's $n = 1$ orbit, \qty{5.29e-11}{m}}{玻尔半径}
  \gterm{Boson}{a particle with whole-number spin}{玻色子}
  % ---- C
  \gterm{Catalyst}{helps a reaction but is not used up}{催化剂}
  \gterm{Cathode}{the negative electrode of a tube}{阴极}
  \gterm{Cathode ray}{a stream of electrons from the cathode}{阴极射线}
  \gterm{Cathode ray tube}{a sealed glass tube with almost no air inside}{阴极射线管}
  \gterm{Cepheid variable}{a pulsing star; its period gives its luminosity}{造父变星}
  \gterm{Chain reaction}{neutrons from fissions cause more fissions}{链式反应}
  \gterm{Charge-to-mass ratio ($q/m$)}{charge divided by mass, in \si{C.kg^{-1}}}{比荷}
  \gterm{Classical physics}{physics before quantum ideas}{经典物理学}
  \gterm{Cloud chamber}{charged particles leave trails of drops in it}{云室}
  \gterm{CNO cycle}{fusion of H to He with C, N, O as catalysts}{碳氮氧循环}
  \gterm{Collider}{a machine where two beams collide head-on}{对撞机}
  \gterm{Collimator}{a slit or hole that makes a narrow beam}{准直器}
  \gterm{Colour charge}{the ``charge'' of the strong force}{色荷}
  \gterm{Confinement}{quarks are never found alone}{夸克禁闭}
  \gterm{Constructive interference}{waves add up in step to make a bright line}{相长干涉}
  \gterm{Continuous spectrum}{every wavelength, with no gaps}{连续光谱}
  \gterm{Control rod}{a rod (boron, cadmium) that absorbs neutrons}{控制棒}
  \gterm{Controlled chain reaction}{exactly one new fission per fission}{受控链式反应}
  \gterm{Coolant}{a fluid that carries heat out of a reactor}{冷却剂}
  \gterm{Cosmic microwave background (CMB)}{faint microwaves from every direction}{宇宙微波背景辐射}
  \gterm{Cosmic rays}{fast particles that arrive from space}{宇宙射线}
  \gterm{Cosmological principle}{on large scales, the Universe is the same}{宇宙学原理}
  \gterm{Count rate}{counts recorded each second (or minute)}{计数率}
  \gterm{Critical mass}{least mass that keeps a chain reaction going}{临界质量}
  \gterm{Crossed fields}{$E$ and $B$ at \ang{90} to each other and the beam}{正交电磁场}
  \gterm{Cyclotron}{particles spiral out in a fixed magnetic field}{回旋加速器}
  % ---- D
  \gterm{Dark energy}{unknown cause of faster expansion (\qty{68}{\percent})}{暗能量}
  \gterm{Dark matter}{unseen matter found by its gravity (\qty{27}{\percent})}{暗物质}
  \gterm{Daughter nucleus}{the new nucleus after a decay}{子核}
  \gterm{de Broglie wavelength}{the wavelength of a moving particle, $\frac{h}{mv}$}{德布罗意波长}
  \gterm{Decay constant ($\lambda$)}{chance per second that one nucleus decays}{衰变常量}
  \gterm{Dee}{a hollow D-shaped electrode of a cyclotron}{D形盒}
  \gterm{Deep inelastic scattering}{fast electrons fired at protons to see inside}{深度非弹性散射}
  \gterm{Detector}{a machine that tracks and measures particles}{探测器}
  \gterm{Deuterium, tritium}{hydrogen-2 and hydrogen-3}{氘，氚}
  \gterm{Diffraction}{a wave spreading out after a gap or an edge}{衍射}
  \gterm{Diffraction grating}{many equal slits that split light into colours}{衍射光栅}
  \gterm{Discharge tube}{a tube of low-pressure gas that glows}{放电管}
  \gterm{Drift tube}{a linac tube that particles coast through}{漂移管}
  % ---- E
  \gterm{Elastic collision}{a collision in which total $K$ stays the same}{弹性碰撞}
  \gterm{Electrode}{a metal plate joined to a power supply}{电极}
  \gterm{Electron}{the light, negative particle in every atom}{电子}
  \gterm{Electron cloud}{a dot picture of where an electron may be}{电子云}
  \gterm{Electron diffraction}{electrons making a wave pattern}{电子衍射}
  \gterm{Electron gun}{makes a narrow beam of electrons}{电子枪}
  \gterm{Electron spin}{makes each electron act like a tiny magnet}{电子自旋}
  \gterm{Electronvolt (eV)}{\qty{1}{eV} $=$ \qty{1.602e-19}{J}}{电子伏特}
  \gterm{Elementary charge ($e$)}{smallest free charge, \qty{1.602e-19}{C}}{元电荷}
  \gterm{e/m tube (fine-beam tube)}{a bulb where coils bend an electron beam}{荷质比测定管}
  \gterm{Emission spectrum}{bright lines at a few exact wavelengths}{发射光谱}
  \gterm{Energy level}{an allowed energy of an electron in an atom}{能级}
  \gterm{Enriched uranium}{uranium with more than \qty{0.7}{\percent} U-235}{浓缩铀}
  \gterm{Event}{one collision in a detector and its products}{事例}
  \gterm{Exchange particle}{a boson that carries a force}{交换粒子}
  \gterm{Excitation}{raising an electron to a higher level}{激发}
  \gterm{Excited state}{any state above the lowest (ground) state}{激发态}
  % ---- F
  \gterm{Fermion}{a particle with half-number spin}{费米子}
  \gterm{Fine structure}{one line that is really several close lines}{精细结构}
  \gterm{Flavour}{the type of a quark (up, down, \ldots)}{味}
  \gterm{Fluorescence}{giving out light as soon as energy hits}{荧光}
  \gterm{Fundamental particle}{a particle with no smaller parts}{基本粒子}
  % ---- G
  \gterm{Galaxy}{billions of stars held together by gravity}{星系}
  \gterm{Gamma ray ($\gamma$)}{a high-energy photon from a nucleus}{γ射线}
  \gterm{Geiger--M\"uller (GM) tube}{a detector that counts ionising radiation}{盖革-米勒计数管}
  \gterm{Generation (of fissions)}{one round of fissions in a chain reaction}{代}
  \gterm{Generation (of particles)}{one of three families of matter particles}{代}
  \gterm{Gluon}{the exchange particle of the strong force}{胶子}
  \gterm{Graviton}{proposed carrier of gravity; never found}{引力子}
  \gterm{Ground state}{the lowest energy level ($n = 1$)}{基态}
  % ---- H
  \gterm{Hadron}{a particle made of quarks}{强子}
  \gterm{Half-life ($t_{\frac{1}{2}}$)}{time for half of the nuclei to decay}{半衰期}
  \gterm{Hertzsprung--Russell (HR) diagram}{stars' luminosity against temperature}{赫罗图}
  \gterm{Higgs boson}{a ripple in the Higgs field, found in 2012}{希格斯玻色子}
  \gterm{Higgs field}{a field that gives most particles their mass}{希格斯场}
  \gterm{Hubble constant ($H_0$)}{slope of a galaxy speed--distance graph}{哈勃常数}
  \gterm{Hubble's law}{recession speed $\propto$ distance: $v = H_0 d$}{哈勃定律}
  \gterm{Hypothesis}{an idea that must be tested}{假说}
  % ---- I
  \gterm{Indirect evidence}{evidence from effects, not direct detection}{间接证据}
  \gterm{Inflation}{a very fast expansion just after the Big Bang}{暴胀}
  \gterm{Ion}{an atom that has lost or gained electrons}{离子}
  \gterm{Ionisation chamber}{a detector that measures a tiny ion current}{电离室}
  \gterm{Ionisation energy}{smallest energy to remove an electron}{电离能}
  \gterm{Ionised}{having lost one or more electrons}{电离的}
  \gterm{Ionising radiation}{radiation that knocks electrons out of atoms}{电离辐射}
  \gterm{Isotope}{same $Z$, different number of neutrons}{同位素}
  % ---- L
  \gterm{Large Hadron Collider (LHC)}{CERN's \qty{27}{km} proton collider}{大型强子对撞机}
  \gterm{Lepton}{a fundamental particle with no strong force}{轻子}
  \gterm{Line spectrum}{only a few sharp wavelengths}{线状光谱}
  \gterm{Linear accelerator (linac)}{speeds particles up along a straight line}{直线加速器}
  \gterm{Luminosity ($L$)}{total power a star gives out, in watts}{光度}
  \gterm{Luminosity class}{I supergiant, III giant, V main sequence}{光度级}
  \gterm{Lyman series}{hydrogen lines ending on $n = 1$ (UV)}{莱曼系}
  % ---- M
  \gterm{Magnitude}{brightness scale: smaller number = brighter}{星等}
  \gterm{Main sequence}{the long stage when a star fuses H into He}{主序}
  \gterm{Maltese cross tube}{a tube where a metal cross casts a shadow}{马耳他十字管}
  \gterm{Mass defect ($\Delta m$)}{mass lost as nucleons join, or in a reaction}{质量亏损}
  \gterm{Mass number ($A$)}{number of protons plus neutrons}{质量数}
  \gterm{Mass--energy equivalence}{mass is a form of energy: $E = mc^2$}{质能等价}
  \gterm{Matter wave}{the wave that goes with a moving particle}{物质波}
  \gterm{Maximum kinetic energy of beta particles ($K_\mathrm{max}$)}{the energy a beta decay releases}{β粒子最大动能}
  \gterm{Megaparsec (Mpc)}{\qty{1}{Mpc} $=$ \qty{3.086e22}{m}}{百万秒差距}
  \gterm{Meson}{a hadron of one quark and one antiquark}{介子}
  \gterm{Model (scientific)}{a simple picture that explains and predicts}{模型}
  \gterm{Moderator}{material that slows neutrons down}{慢化剂}
  \gterm{Momentum ($p$)}{mass times velocity, $p = mv$}{动量}
  \gterm{Muon}{a heavy copy of the electron (207 times)}{μ子}
  % ---- N
  \gterm{Natural logarithm (ln)}{the function that undoes $e^x$}{自然对数}
  \gterm{Nebula}{a cloud of gas and dust in space}{星云}
  \gterm{Neutral (atom)}{an atom with all its electrons: no net charge}{中性原子}
  \gterm{Neutrino}{a neutral lepton with a tiny mass}{中微子}
  \gterm{Neutron}{a neutral particle in the nucleus}{中子}
  \gterm{Neutron star}{a dead core about \qty{20}{km} across, of neutrons}{中子星}
  \gterm{Nuclear fission}{a heavy nucleus splits into two smaller ones}{核裂变}
  \gterm{Nuclear fusion}{light nuclei join to make a heavier one}{核聚变}
  \gterm{Nuclear model}{Rutherford's atom, with a tiny, dense nucleus}{核式结构模型}
  \gterm{Nucleon}{a proton or a neutron}{核子}
  \gterm{Nucleosynthesis}{the making of new nuclei}{核合成}
  \gterm{Nucleus}{the tiny, dense, positive centre of an atom}{原子核}
  % ---- O
  \gterm{Orbital}{a region where an electron is likely to be}{原子轨道}
  % ---- P
  \gterm{Paddle wheel tube}{a tube with a light wheel turned by the beam}{叶轮管}
  \gterm{Pair production}{a photon becomes a particle--antiparticle pair}{对产生}
  \gterm{Paraffin wax}{a solid rich in hydrogen (protons)}{石蜡}
  \gterm{Parent nucleus}{the nucleus before a decay}{母核}
  \gterm{Parsec (pc)}{a distance: \qty{1}{pc} $=$ 3.26 light-years}{秒差距}
  \gterm{Particle accelerator}{a machine that makes charged particles fast}{粒子加速器}
  \gterm{Particle zoo}{the many particles found from 1950 to 1970}{粒子动物园}
  \gterm{Paschen series}{hydrogen lines ending on $n = 3$ (IR)}{帕邢系}
  \gterm{Pauli exclusion principle}{no two electrons share all quantum numbers}{泡利不相容原理}
  \gterm{Penetrating power}{how far radiation travels through a material}{穿透能力}
  \gterm{Permittivity of free space ($\varepsilon_0$)}{the constant in Coulomb's law}{真空电容率}
  \gterm{Photon}{a packet of light energy, $E = hf$}{光子}
  \gterm{Pion}{the lightest meson}{π介子}
  \gterm{Planetary nebula}{gas blown off by a dying Sun-like star}{行星状星云}
  \gterm{Plasma}{a gas so hot that atoms lose their electrons}{等离子体}
  \gterm{Plum pudding model}{Thomson's atom: electrons in a $+$ ball}{枣糕模型}
  \gterm{Point-like}{so small that it acts like a dot}{点状的}
  \gterm{Positron}{the electron's antiparticle, charge $+e$}{正电子}
  \gterm{Postulate}{a rule assumed without proof}{公设}
  \gterm{Probability}{how likely: from 0 (never) to 1 (certain)}{概率}
  \gterm{Proton}{a positive particle in the nucleus (uud)}{质子}
  \gterm{Proton--proton chain}{fusion of 4 protons into He-4 (Sun-like stars)}{质子-质子链}
  \gterm{Protostar}{a shrinking gas clump before fusion starts}{原恒星}
  % ---- Q
  \gterm{Quantised}{only certain separate values are allowed}{量子化的}
  \gterm{Quantum number}{a whole number such as $n$ that labels a state}{量子数}
  \gterm{Quark}{a fundamental particle inside hadrons}{夸克}
  % ---- R
  \gterm{Radioactive decay}{an unstable nucleus changes by itself}{放射性衰变}
  \gterm{Radioisotope}{an isotope that is radioactive}{放射性同位素}
  \gterm{Random}{cannot be predicted for one nucleus}{随机的}
  \gterm{Recession velocity ($v$)}{speed at which a galaxy moves away}{退行速度}
  \gterm{Recoil}{to move backwards after a hit or an emission}{反冲}
  \gterm{Red giant}{a big, cool, bright star late in its life}{红巨星}
  \gterm{Red supergiant}{a huge, cool star that fuses up to iron}{红超巨星}
  \gterm{Redshift}{lines moved to longer wavelengths}{红移}
  \gterm{Relativistic momentum ($p_\mathrm{v}$)}{momentum of a particle near $c$}{相对论动量}
  \gterm{Reliable (reliability)}{gives the same results when repeated}{可靠性}
  \gterm{RF cavity}{where an electric field pushes the beam}{射频腔}
  \gterm{Rydberg constant ($R$)}{\qty{1.097e7}{m^{-1}}, used for hydrogen}{里德伯常量}
  % ---- S
  \gterm{Scintillation}{a tiny flash when a particle hits a screen}{闪烁}
  \gterm{Series limit}{shortest wavelength of a series ($n_\mathrm{i} = \infty$)}{线系极限}
  \gterm{Solar mass ($M_\odot$)}{the Sun's mass, used as a unit}{太阳质量}
  \gterm{Spectral class}{O B A F G K M, from hot to cool}{光谱型}
  \gterm{Spectral line}{a bright or dark line at one wavelength}{谱线}
  \gterm{Spectrograph}{spreads light into a spectrum and records it}{摄谱仪}
  \gterm{Spectrometer}{measures the angles of spectral lines}{分光计}
  \gterm{Spectrum}{light spread out into its wavelengths}{光谱}
  \gterm{Spin}{a built-in amount of turning of a particle}{自旋}
  \gterm{Spontaneous}{happening by itself, with no trigger}{自发的}
  \gterm{Stable}{lasts a long time without changing}{稳定的}
  \gterm{Standard Model}{theory of fundamental particles and forces}{标准模型}
  \gterm{Standing wave}{a wave pattern with fixed nodes}{驻波}
  \gterm{Stark effect}{a line splits in an electric field}{斯塔克效应}
  \gterm{Stationary state}{an allowed orbit, with no radiation}{定态}
  \gterm{Stellar equilibrium}{inward gravity balances outward pressure}{流体静力平衡}
  \gterm{Strong force}{holds quarks together; carried by gluons}{强相互作用}
  \gterm{Strong nuclear force}{short-range attraction between nucleons}{核力}
  \gterm{Subatomic particle}{a particle smaller than an atom}{亚原子粒子}
  \gterm{Subclass}{0 (hottest) to 9 (coolest) in a spectral class}{光谱次型}
  \gterm{Subcritical}{too little fuel: the chain reaction dies out}{次临界}
  \gterm{Supercritical}{more than critical: fissions grow fast}{超临界}
  \gterm{Supernova}{the explosion of a massive star}{超新星}
  \gterm{Synchrocyclotron}{a cyclotron whose frequency falls near $c$}{同步回旋加速器}
  \gterm{Synchrotron}{a ring whose $B$ rises as the beam speeds up}{同步加速器}
  \gterm{Systematic error}{an error that shifts every result one way}{系统误差}
  % ---- T
  \gterm{Terminal velocity}{steady falling speed: drag equals weight}{收尾速度}
  \gterm{Thermal (slow) neutron}{a slow neutron that U-235 captures easily}{热中子}
  \gterm{Thermionic emission}{a hot metal giving off electrons}{热电子发射}
  \gterm{Transition}{a jump between two energy levels}{跃迁}
  \gterm{Transmutation}{a nuclear change to another element}{核嬗变}
  \gterm{Triple-alpha process}{three helium-4 nuclei fuse into carbon-12}{3α过程}
  % ---- U
  \gterm{Uncertainty principle}{position and momentum are not both exact}{不确定性关系}
  \gterm{Uncontrolled chain reaction}{more than one new fission per fission}{不受控链式反应}
  \gterm{Uniform field}{the same strength and direction everywhere}{匀强场}
  % ---- V
  \gterm{Valid (validity)}{tests what it claims to test}{有效性}
  \gterm{Velocity selector}{crossed fields that pass only one speed}{速度选择器}
  % ---- W
  \gterm{W and Z bosons}{the exchange particles of the weak force}{W和Z玻色子}
  \gterm{Wave equation (Schr\"odinger equation)}{how the electron's wave behaves in an atom}{薛定谔方程}
  \gterm{Wave function ($\psi$)}{the electron's wave; $\psi^2$ gives the probability}{波函数}
  \gterm{Wave--particle duality}{light and matter act as waves and particles}{波粒二象性}
  \gterm{Weak force}{changes quark flavour, as in beta decay}{弱相互作用}
  \gterm{White dwarf}{small, hot, dim core left by a Sun-like star}{白矮星}
  \gterm{Wien's displacement law}{$\lambda_\mathrm{max} = \frac{b}{T}$: hotter means shorter}{维恩位移定律}
  % ---- Z
  \gterm{Zeeman effect}{a line splits in a magnetic field}{塞曼效应}
  \gterm{Zero error}{a scale that does not read zero at zero}{零点误差}
\end{glossarytable}
